% !TeX program = lualatex
% ================================================================
% French Mathematical Dictionary
%
% Generated by the server using the following settings:
%
%   language            French (fra)
%   text direction      left to right
%   font                Noto Serif
%   fallback font       Noto Serif
%   built from          latex/dictionary/mathematicalDictionary.tex
%   tier 3 vocabulary   green   RGB 0 128 0
%   English text        black   RGB 0 0 0
%   French text         purple  RGB 112 48 160
%   layout macros       version 3
%
% Please don't edit any information in this comment block - the server
% compares it against the current settings and rebuilds this file when
% they differ, so changes here would be overwritten. However, feel free
% to make updates to the LaTeX below if you want to edit it.
% ================================================================
%
% ---- What is safe to edit in this file ----------------------------------
%
% Safe:     the French word and the French definition - the third
%           and fourth arguments of each entry. That is what this file is for:
%           if a translation reads wrongly to somebody who speaks the language,
%           correct it here and every sheet that uses the word follows.
%
% Not safe: the English word and the English definition - the first and second
%           arguments. The first is how an entry is matched to the shared
%           dictionary; the second is the wording the translation was made
%           from, and is how the server knows the translation is still current.
%           Changing an English definition here does not change the shared
%           dictionary - it only makes this entry look up to date when it is
%           not. Edit the English in the shared dictionary instead, and this
%           file will be brought back into step on the next run.
%
% Entries are added and removed by the server, following the shared dictionary.
% An entry the server cannot read is reported rather than guessed at, so if a
% run complains about this file, look for a missing brace.
% -------------------------------------------------------------------------

\documentclass[11pt]{article}
\usepackage[a4paper,margin=18mm]{geometry}
% Characters this language's font has not got are borrowed from another,
% rather than being silently dropped - see docs/translated-sheet-latex.md
\directlua{%
  if luaotfload and luaotfload.add_fallback then
    luaotfload.add_fallback("eallatinfallback",
      { "Noto Serif:mode=harf;" })
  else
    tex.error("this luaotfload is too old to register a fallback font")
  end
}
\usepackage[english]{babel}
\babelprovide[import]{french}
\babelfont[french]{rm}[Renderer=Harfbuzz, BoldFont={Noto Serif}, ItalicFont={Noto Serif}, BoldItalicFont={Noto Serif}, RawFeature={fallback=eallatinfallback}]{Noto Serif}
\babelfont[french]{sf}[Renderer=Harfbuzz, BoldFont={Noto Serif}, ItalicFont={Noto Serif}, BoldItalicFont={Noto Serif}, RawFeature={fallback=eallatinfallback}]{Noto Serif}
\newcommand{\ealtext}[1]{\foreignlanguage{french}{#1}}
\newcommand{\ealtextblock}[1]{{\raggedright\foreignlanguage{french}{#1}\par}}
% ================================================================
% EAL layout helpers
% ================================================================
\usepackage{xcolor}

\definecolor{ealkeycolour}{RGB}{0,128,0}
\definecolor{ealenglish}{RGB}{0,0,0}
\definecolor{ealtwocolour}{RGB}{112,48,160}

\newsavebox{\ealboxen}
\newsavebox{\ealboxtr}
\newlength{\ealglosswd}
\newlength{\ealglossraise}
\setlength{\ealglossraise}{1.15em}

% a tier 3 word, in English and in translation alike
\newcommand{\ealkey}[1]{\textcolor{ealkeycolour}{#1}}
\newcommand{\ealkeytr}[1]{\textcolor{ealkeycolour}{\ealtext{#1}}}

% One translated word above one English word. Built from kernel
% primitives, never a tabular, which would break wherever the sheet
% loads array, booktabs or tabularx. The box is as wide as the wider of
% the two, so a long translation pushes its neighbours aside instead of
% overlapping them, and it claims its full height so a table row or a
% TikZ node makes room for it.
\newcommand{\ealgloss}[2]{%
  \sbox{\ealboxen}{\ealkey{#1}}%
  \sbox{\ealboxtr}{\small\textcolor{ealtwocolour}{\ealtext{#2}}}%
  \setlength{\ealglosswd}{\wd\ealboxen}%
  \ifdim\wd\ealboxtr>\ealglosswd\setlength{\ealglosswd}{\wd\ealboxtr}\fi
  \leavevmode
  \makebox[\ealglosswd][c]{%
    \makebox[0pt][c]{\usebox{\ealboxen}}%
    \makebox[0pt][c]{\raisebox{\ealglossraise}{\usebox{\ealboxtr}}}%
  }%
}

% opens up line spacing around a block of glosses, so the raised
% translations sit clear of the line above
\newenvironment{ealglossed}{\par\linespread{2.1}\selectfont\ignorespaces}{\par}

% a whole translated sentence above its English counterpart. block level,
% so it does not belong inside a tabular cell
\newcommand{\ealpara}[2]{%
  \par\addvspace{0.5\baselineskip}%
  {\color{ealtwocolour}\ealtextblock{#1}}%
  \penalty200
  {\color{ealenglish}#2\par}%
  \addvspace{0.5\baselineskip}%
}

% ================================================================
% Dictionary layout
% ================================================================
\usepackage{multicol}

\newlength{\ealdictgap}
\setlength{\ealdictgap}{1.2mm}

% english word / english meaning / translated word / translated meaning
\newcommand{\dictentrytr}[4]{%
  \par\addvspace{\ealdictgap}%
  \noindent
  {\bfseries\ealkey{#1}}\quad{\small #2}\par
  \nopagebreak
  {\color{ealtwocolour}\ealtext{{\bfseries #3}~\textemdash{} {\small #4}}\par}%
  \par
}

\pagestyle{empty}
\setlength{\parindent}{0pt}

\begin{document}
{\large\bfseries French Mathematical Dictionary}\par\medskip

\dictentrytr{absolute value}{the distance of a number from zero, ignoring whether it is positive or negative.}{valeur absolue}{la distance d'un nombre à zéro, sans tenir compte de son signe positif ou négatif.}
\dictentrytr{addition}{combining two or more numbers to find their total.}{addition}{combiner deux nombres ou plus pour trouver leur total.}
\dictentrytr{adjacent}{next to, and sharing a side or a corner}{adjacent}{à côté, et partageant un côté ou un sommet}
\dictentrytr{altitude}{height above sea level or another fixed reference level.}{altitude}{hauteur au-dessus du niveau de la mer ou d'un autre niveau de référence fixe.}
\dictentrytr{angle}{the amount of turn between two straight lines that meet at a point}{angle}{la quantité de rotation entre deux droites qui se rencontrent en un point}
\dictentrytr{annular}{Describing a shape made by the space between two concentric circles.}{annulaire}{Décrivant une forme constituée par l'espace entre deux cercles concentriques.}
\dictentrytr{apex angle}{the angle at the point of a triangle opposite its base.}{angle au sommet}{l'angle situé au sommet d'un triangle, opposé à sa base.}
\dictentrytr{arc}{part of the curved edge of a circle}{arc}{partie du bord courbe d'un cercle}
\dictentrytr{arccos}{The operation that finds an angle when its cosine is known.}{arccos}{L'opération qui trouve un angle lorsque son cosinus est connu.}
\dictentrytr{arctan}{The operation that finds an angle when its tangent is known.}{arctan}{L'opération qui trouve un angle lorsque sa tangente est connue.}
\dictentrytr{area}{the amount of space inside a flat shape}{aire}{la quantité d'espace à l'intérieur d'une forme plane}
\dictentrytr{area model}{a rectangle diagram split into smaller rectangles, used to show how to multiply fractions.}{modèle d'aire}{un schéma en rectangle divisé en petits rectangles, utilisé pour montrer comment multiplier des fractions.}
\dictentrytr{arithmetic}{the branch of mathematics that deals with adding, subtracting, multiplying and dividing numbers.}{arithmétique}{la branche des mathématiques qui traite de l'addition, de la soustraction, de la multiplication et de la division des nombres.}
\dictentrytr{ascending order}{arranged from smallest to largest}{ordre croissant}{rangé du plus petit au plus grand}
\dictentrytr{average}{a single value chosen to stand for a whole set of data}{moyenne}{une valeur unique choisie pour représenter tout un ensemble de données}
\dictentrytr{average speed}{total distance travelled divided by total time taken.}{vitesse moyenne}{distance totale parcourue divisée par le temps total mis.}
\dictentrytr{axis}{one of the lines a graph is measured against}{axe}{l'une des droites par rapport auxquelles un graphique est mesuré}
\dictentrytr{bar model}{a diagram made of rectangles that shows the parts and total of an equation.}{modèle en barres}{un schéma fait de rectangles qui montre les parties et le total d'une équation.}
\dictentrytr{base}{the side of a shape chosen as the bottom for working out its area.}{base}{le côté d'une forme choisi comme fond pour calculer son aire.}
\dictentrytr{basis point}{a unit equal to one hundredth of one percent, used for small changes in rates.}{point de base}{une unité égale à un centième d'un pour cent, utilisée pour de petites variations de taux.}
\dictentrytr{bearing}{a direction given as an angle clockwise from north, written with three figures}{azimut}{une direction donnée sous forme d'un angle mesuré dans le sens horaire depuis le nord, écrite avec trois chiffres}
\dictentrytr{bimodal}{having two values that appear most often in a list of numbers.}{bimodal}{ayant deux valeurs qui apparaissent le plus souvent dans une liste de nombres.}
\dictentrytr{bisector}{a line that cuts an angle or another line exactly in half}{bissectrice}{une droite qui coupe un angle ou une autre droite exactement en deux}
\dictentrytr{bodmas}{a rule that says which operation to do first in a calculation.}{règle des priorités opératoires}{une règle qui indique quelle opération effectuer en premier dans un calcul.}
\dictentrytr{bond}{a type of loan to a government or company that pays regular money and is paid back later.}{obligation}{un type de prêt à un gouvernement ou à une entreprise qui verse de l'argent régulièrement et est remboursé plus tard.}
\dictentrytr{box method}{a way of multiplying or expanding by splitting a rectangle into smaller parts.}{méthode des boîtes}{une façon de multiplier ou de développer en divisant un rectangle en parties plus petites.}
\dictentrytr{bracket}{one of a pair of curved symbols, ( ), used to group numbers or letters.}{parenthèse}{l'un d'une paire de symboles courbes, ( ), utilisés pour regrouper des nombres ou des lettres.}
\dictentrytr{brackets}{A pair of signs used in mathematics to group together a part of an expression.}{parenthèses}{Une paire de signes utilisée en mathématiques pour regrouper une partie d'une expression.}
\dictentrytr{cancel}{to remove a shared number from the top and bottom of a fraction before multiplying.}{simplifier}{retirer un nombre commun en haut et en bas d'une fraction avant de multiplier.}
\dictentrytr{cardinal directions}{The four main compass points: North, South, East and West.}{points cardinaux}{Les quatre points principaux de la boussole : Nord, Sud, Est et Ouest.}
\dictentrytr{category}{a group of non-numerical items, for example a colour, used as data.}{catégorie}{un groupe d'éléments non numériques, par exemple une couleur, utilisé comme données.}
\dictentrytr{centre}{the fixed middle point of a circle}{centre}{le point milieu fixe d'un cercle}
\dictentrytr{chord}{a straight line joining two points on the edge of a circle}{corde}{une droite joignant deux points du bord d'un cercle}
\dictentrytr{circle}{a round flat shape whose edge is always the same distance from its centre}{cercle}{une forme plate ronde dont le bord est toujours à la même distance de son centre}
\dictentrytr{circular}{round in shape, with every point on its edge the same distance from the centre.}{circulaire}{de forme ronde, avec chaque point de son bord à la même distance du centre.}
\dictentrytr{circumference}{the distance all the way round a circle}{circonférence}{la distance sur tout le tour d'un cercle}
\dictentrytr{coefficient}{the number in front of a letter, showing how many of it there are}{coefficient}{le nombre devant une lettre, indiquant combien il y en a}
\dictentrytr{column vector}{a vector written vertically with one number above another.}{vecteur colonne}{un vecteur écrit verticalement avec un nombre au-dessus d'un autre.}
\dictentrytr{common denominator}{a bottom number shared by two or more fractions, used for adding or comparing them.}{dénominateur commun}{un nombre du bas partagé par deux fractions ou plus, utilisé pour les additionner ou les comparer.}
\dictentrytr{common difference}{the fixed number added to each term to get the next term}{raison}{le nombre fixe ajouté à chaque terme pour obtenir le terme suivant}
\dictentrytr{common factor}{a number that divides exactly into two or more other numbers.}{facteur commun}{un nombre qui divise exactement deux autres nombres ou plus.}
\dictentrytr{common ratio}{the fixed number used to create each next term in a geometric sequence}{raison}{le nombre fixe utilisé pour créer chaque terme suivant dans une suite géométrique}
\dictentrytr{congruent}{exactly the same shape and size}{congru}{exactement la même forme et la même taille}
\dictentrytr{correlation}{how closely two sets of data are related}{corrélation}{dans quelle mesure deux ensembles de données sont liés}
\dictentrytr{cumulative}{adding up as you go along, so each total includes the ones before}{cumulatif}{s'additionnant au fur et à mesure, de sorte que chaque total inclut les précédents}
\dictentrytr{data}{the set of values collected}{données}{l'ensemble des valeurs collectées}
\dictentrytr{decimal}{a number written with a point, where the digits after the point are parts of a whole}{décimal}{un nombre écrit avec une virgule, où les chiffres après la virgule sont des parties d'un tout}
\dictentrytr{denominator}{the number below the line in a fraction, showing how many equal parts the whole is split into}{dénominateur}{le nombre sous la barre d'une fraction, indiquant en combien de parties égales le tout est partagé}
\dictentrytr{diameter}{a straight line across a circle through its centre}{diamètre}{une droite traversant un cercle par son centre}
\dictentrytr{difference}{the answer when one number is subtracted from another}{différence}{le résultat lorsque l'on soustrait un nombre à un autre}
\dictentrytr{digit}{a single figure from 0 to 9 that a number is written with}{chiffre}{un seul chiffre de 0 à 9 avec lequel un nombre s'écrit}
\dictentrytr{edge}{a straight line where two faces of a solid meet}{arête}{une droite où deux faces d'un solide se rencontrent}
\dictentrytr{equivalent}{having the same value, even though it is written differently}{équivalent}{ayant la même valeur, même si c'est écrit différemment}
\dictentrytr{estimate}{a sensible answer worked out quickly from rounded numbers}{estimation}{une réponse raisonnable calculée rapidement à partir de nombres arrondis}
\dictentrytr{expand}{to multiply out brackets so the expression has no brackets left}{développer}{multiplier pour enlever les parenthèses afin que l'expression n'ait plus de parenthèses}
\dictentrytr{expression}{a group of numbers and letters joined by operations, with no equals sign}{expression}{un groupe de nombres et de lettres reliés par des opérations, sans signe égal}
\dictentrytr{face}{a flat surface of a solid shape}{face}{une surface plane d'une forme solide}
\dictentrytr{factor}{a whole number that divides exactly into another number}{facteur}{un nombre entier qui divise exactement un autre nombre}
\dictentrytr{factorise}{to write an expression as a product, usually by putting brackets back in}{factoriser}{écrire une expression sous forme de produit, généralement en remettant des parenthèses}
\dictentrytr{formula}{a rule written with letters, showing how to work one quantity out from others}{formule}{une règle écrite avec des lettres, montrant comment calculer une quantité à partir d'autres}
\dictentrytr{fraction}{a number written as one whole split into equal parts, with a number of those parts taken}{fraction}{un nombre écrit comme un tout divisé en parties égales, avec un certain nombre de ces parties prises}
\dictentrytr{frequency}{how many times something happens}{effectif}{le nombre de fois qu'une chose se produit}
\dictentrytr{gradient}{how steep a line is, found by how far it rises for each step across}{pente}{à quel point une droite est pentue, obtenu par la distance qu'elle monte pour chaque pas horizontal}
\dictentrytr{highest common factor}{the largest whole number that divides exactly into two or more numbers}{plus grand commun diviseur}{le plus grand nombre entier qui divise exactement deux nombres ou plus}
\dictentrytr{hypotenuse}{the longest side of a right-angled triangle, opposite the right angle}{hypoténuse}{le côté le plus long d'un triangle rectangle, opposé à l'angle droit}
\dictentrytr{identity}{a statement that is true for every value of the letters in it}{identité}{une affirmation qui est vraie pour toutes les valeurs des lettres qu'elle contient}
\dictentrytr{index}{the small raised number showing what power something is raised to}{exposant}{le petit nombre en exposant indiquant à quelle puissance quelque chose est élevé}
\dictentrytr{inequality}{a statement that one quantity is larger or smaller than another}{inégalité}{une affirmation qu'une quantité est plus grande ou plus petite qu'une autre}
\dictentrytr{integer}{a whole number, which may be positive, negative or zero}{entier}{un nombre entier, qui peut être positif, négatif ou zéro}
\dictentrytr{intercept}{the point where a graph crosses an axis}{point d'intersection avec un axe}{le point où un graphique coupe un axe}
\dictentrytr{linear}{making a straight line when drawn as a graph}{linéaire}{qui forme une ligne droite lorsqu'on le dessine sous forme de graphique}
\dictentrytr{lowest common multiple}{the smallest number that two or more numbers all divide into exactly}{plus petit commun multiple}{le plus petit nombre que deux nombres ou plus divisent tous exactement}
\dictentrytr{mean}{the average found by adding all the values and dividing by how many there are}{moyenne}{la moyenne obtenue en additionnant toutes les valeurs et en divisant par leur nombre}
\dictentrytr{median}{the middle value when the data is put in order}{médiane}{la valeur du milieu lorsque les données sont rangées dans l'ordre}
\dictentrytr{mode}{the value that appears most often}{mode}{la valeur qui apparaît le plus souvent}
\dictentrytr{multiple}{a number you get by multiplying another number by a whole number}{multiple}{un nombre que l'on obtient en multipliant un autre nombre par un nombre entier}
\dictentrytr{multiply}{to add a number to itself a given number of times}{multiplier}{ajouter un nombre à lui-même un nombre donné de fois}
\dictentrytr{nth term}{a rule that gives any number in a sequence from its position}{terme de rang n}{une règle qui donne n'importe quel nombre d'une suite à partir de sa position}
\dictentrytr{numerator}{the number above the line in a fraction, showing how many equal parts are taken}{numérateur}{le nombre au-dessus de la barre d'une fraction, indiquant combien de parties égales sont prises}
\dictentrytr{outlier}{a value far away from all the others}{valeur aberrante}{une valeur très éloignée de toutes les autres}
\dictentrytr{parallel}{always the same distance apart and never meeting}{parallèle}{toujours à la même distance l'une de l'autre et ne se rencontrant jamais}
\dictentrytr{percentage}{a number of parts out of a hundred}{pourcentage}{un nombre de parties sur cent}
\dictentrytr{perimeter}{the total distance around the edge of a shape}{périmètre}{la distance totale autour du bord d'une forme}
\dictentrytr{perpendicular}{meeting at a right angle}{perpendiculaire}{se rencontrant à angle droit}
\dictentrytr{place value}{the value a digit has because of where it sits in a number}{valeur de position}{la valeur qu'un chiffre a en raison de sa position dans un nombre}
\dictentrytr{polygon}{a flat shape with straight sides}{polygone}{une forme plane à côtés droits}
\dictentrytr{power}{how many times a number is multiplied by itself}{puissance}{combien de fois un nombre est multiplié par lui-même}
\dictentrytr{prime}{a whole number greater than 1 that divides only by itself and 1}{premier}{un nombre entier supérieur à 1 qui n'est divisible que par lui-même et par 1}
\dictentrytr{probability}{how likely something is, given as a number between 0 and 1}{probabilité}{la probabilité qu'une chose se produise, donnée sous la forme d'un nombre entre 0 et 1}
\dictentrytr{product}{the answer when two or more numbers are multiplied}{produit}{le résultat lorsque deux nombres ou plus sont multipliés}
\dictentrytr{proportion}{a relationship where two amounts change by the same factor}{proportion}{une relation où deux quantités changent par le même facteur}
\dictentrytr{quadratic}{an expression or equation whose highest power is a square}{quadratique}{une expression ou une équation dont la puissance la plus élevée est un carré}
\dictentrytr{quartile}{a value marking one quarter of the way through ordered data}{quartile}{une valeur marquant un quart dans des données ordonnées}
\dictentrytr{quotient}{the answer when one number is divided by another}{quotient}{le résultat lorsque l'on divise un nombre par un autre}
\dictentrytr{radius}{the distance from the centre of a circle to its edge}{rayon}{la distance du centre d'un cercle à son bord}
\dictentrytr{range}{the difference between the largest and smallest values}{étendue}{la différence entre la plus grande et la plus petite des valeurs}
\dictentrytr{ratio}{a way of comparing two or more amounts, showing how much of one there is for each of the other}{rapport}{une façon de comparer deux quantités ou plus, montrant combien il y a de l'une pour chacune des autres}
\dictentrytr{rearrange}{to change the order of a formula so a different letter is on its own}{réarranger}{changer l'ordre d'une formule pour qu'une autre lettre soit seule}
\dictentrytr{reciprocal}{the number you multiply by to get 1, found by turning a fraction upside down}{inverse}{le nombre par lequel on multiplie pour obtenir 1, que l'on trouve en retournant une fraction}
\dictentrytr{recur}{to repeat the same digits forever after the decimal point}{se répéter}{répéter indéfiniment les mêmes chiffres après la virgule}
\dictentrytr{root}{a number that gives another number when multiplied by itself a given number of times}{racine}{un nombre qui donne un autre nombre lorsqu'on le multiplie par lui-même un nombre donné de fois}
\dictentrytr{round}{to replace a number with a nearby simpler one}{arrondir}{remplacer un nombre par un nombre plus simple proche}
\dictentrytr{sample}{a smaller group chosen to stand for a larger one}{échantillon}{un groupe plus petit choisi pour représenter un groupe plus grand}
\dictentrytr{scale factor}{the number every length is multiplied by when a shape is enlarged}{facteur d'échelle}{le nombre par lequel chaque longueur est multipliée lorsqu'une figure est agrandie}
\dictentrytr{sector}{a slice of a circle between two radii}{secteur}{une portion de cercle entre deux rayons}
\dictentrytr{segment}{the part of a circle cut off by a chord}{segment circulaire}{la partie d'un cercle découpée par une corde}
\dictentrytr{sequence}{a list of numbers that follows a rule}{suite}{une liste de nombres qui suit une règle}
\dictentrytr{similar}{the same shape but a different size}{semblable}{la même forme mais de taille différente}
\dictentrytr{simplify}{to write something in its smallest or plainest form without changing its value}{simplifier}{écrire quelque chose sous sa forme la plus petite ou la plus simple sans changer sa valeur}
\dictentrytr{substitute}{to put a number in place of a letter}{substituer}{mettre un nombre à la place d'une lettre}
\dictentrytr{sum}{the answer when numbers are added together}{somme}{le résultat lorsque des nombres sont additionnés}
\dictentrytr{surd}{a root that cannot be written exactly as a fraction, so it is left in root form}{racine irrationnelle}{une racine qui ne peut pas s'écrire exactement sous forme de fraction, donc on la laisse sous forme de racine}
\dictentrytr{tangent}{a line that touches a curve at exactly one point}{tangente}{une droite qui touche une courbe en exactement un point}
\dictentrytr{term}{one part of an expression or one number in a sequence}{terme}{une partie d'une expression ou un nombre dans une suite}
\dictentrytr{transformation}{a change to a shape's position, size or orientation}{transformation}{un changement de position, de taille ou d'orientation d'une forme}
\dictentrytr{unitary}{a method that first finds the value of one part or one item}{unitaire}{une méthode qui trouve d'abord la valeur d'une part ou d'un élément}
\dictentrytr{variable}{a letter standing for a quantity that can change}{variable}{une lettre qui représente une quantité pouvant changer}
\dictentrytr{vector}{a quantity with both size and direction}{vecteur}{une quantité avec à la fois une taille et une direction}
\dictentrytr{vertex}{a corner where two or more edges meet}{sommet}{un coin où deux arêtes ou plus se rencontrent}
\dictentrytr{volume}{the amount of space inside a solid shape}{volume}{la quantité d'espace à l'intérieur d'une forme solide}

\end{document}
